% Round 21 / Hodge branch. New candidate; independent review required.
\section{Mixed surface--hyperplane smoothings with full local Green theory}
\label{r21mx:section}

\paragraph{Scope and relation to the frozen families.}
We give an actual fourfold smoothing whose central factors lie in
different line bundles and are not ample divisors in the ambient
fivefold.  Its component and double-stratum geometry closes both the
middle numerical gate and the native curve-comparison gate.  A separate
weight calculation produces members with no potential good reduction.
This supplies a newly verified local presentation in this archive; it
does not claim disjointness from every earlier abstract motivic
criterion.  No new global number-field coefficient family or complete
height family is asserted without the required witnesses at all places.

The ordered local resolution calculation is the one used in the frozen
HP18/NF18/PG18 sources.  We explain why its geometric part applies to
these different factor line bundles.  We do not apply PG18's numerical
theorem outside its common-ample-power hypotheses.  Instead the
numerical argument uses the full component intersection forms and the
reviewed PU18/FP21 criterion.

\subsection{The supplied native datum and the actual resolution}

Let $R$ be the ring of integers of a finite $K/\mathbf Q_p$, with its
fixed absolute value and uniformizer $\pi$, and let $k$ be its finite
residue field.  Let $\mathcal S/R$ be smooth projective with
geometrically integral surface fibers.  Choose a relatively very ample
line bundle $\mathcal A$ and a section $f_0\in H^0(\mathcal S,\mathcal A)$
whose special zero locus $C\subset S=\mathcal S_k$ is a smooth
geometrically integral curve.  Fix $d\ge4$ and linear forms
$l_1,\ldots,l_d$ on $\mathbf P_R^3$ whose special hyperplanes are in
geometric general position: any at most three meet in the expected
projective linear space and any four have empty intersection.

Put $\mathcal Y=\mathcal S\times_R\mathbf P_R^3$ and
$\mathcal L=\mathcal A\boxtimes\mathcal O(d)$.  Let
$Q\in H^0(\mathcal Y,\mathcal L)$.  On $Y=\mathcal Y_k$ set
\[
 H_0=C\times\mathbf P_k^3,\qquad
 H_j=S\times V(\bar l_j)\quad(1\le j\le d),
 \qquad H_I=\bigcap_{i\in I}H_i.
 \tag{MX21.1}\label{r21mx:central}
\]
Assume that $\bar Q$ cuts every nonempty $H_I$ transversely, that
its positive-dimensional intersections with these strata are
geometrically integral, and that all zero-dimensional intersections
are finite etale.  Assume also that
\[
 X=V\left(f_0\prod_{j=1}^d l_j+\pi Q\right)
                         \subset\mathcal Y_K
 \tag{MX21.2}\label{r21mx:generic}
\]
is smooth and geometrically integral.  These are explicit geometric
hypotheses on the chosen sections.  Their simultaneous realizability
after a finite unramified extension is proved below.

The nonempty central strata are exactly
\[
 \begin{array}{ll}
 H_I\simeq S\times\mathbf P^{3-|I|},
       &0\notin I,\quad1\le |I|\le3,\\[2pt]
 H_I\simeq C\times\mathbf P^{4-|I|},
       &0\in I,\quad1\le |I|\le4.
 \end{array}
 \tag{MX21.3}\label{r21mx:strata}
\]
They are smooth and geometrically integral of dimension $5-|I|$.
In particular triple and quadruple strata are nonempty.  There are
no quintuple strata.

Blow up the smooth ambient $\mathcal Y$ successively in the strict
transforms of the special subvarieties $H_0,H_1,\ldots,H_{d-1}$,
and take the strict transform $\mathcal X$ of the flat closure of
\eqref{r21mx:generic}.

\begin{proposition}[Actual ordered-resolution geometry]
\label{r21mx:resolution}
The scheme $\mathcal X/R$ is regular, flat and projective, with
generic fiber $X$ and reduced strict semistable special fiber
$\bigcup_{i=0}^dD_i$.  No additional vertical irreducible component
occurs.  For a nonempty index set $I$, its actual closed stratum
$D_I=\bigcap_{i\in I}D_i$ is the iterated smooth blowup of $H_I$ in
the reverse-order centers
\[
          Z_{I,j}=H_I\cap H_j\cap V(\bar Q),
                    \qquad j=d,d-1,\ldots,\max(I)+1.
 \tag{MX21.4}\label{r21mx:centers}
\]
Empty centers are omitted.  Every occurring center has codimension
two in $H_I$, and its strict transform in the indicated sequence is
isomorphic to the original center.  Thus each component is obtained
from $C\times\mathbf P^3$ or $S\times\mathbf P^2$ by smooth surface
blowups, and each double stratum from $C\times\mathbf P^2$ or
$S\times\mathbf P^1$ by smooth curve blowups.
\end{proposition}

\begin{proof}
At a special point on $s$ central factors the local defining functions
of these factors are part of a regular relative parameter system,
by \eqref{r21mx:strata}.  If $\bar Q$ is a unit, the equation is
$x_1\cdots x_s=\pi u$ with $u$ a unit.  It is already strict
semistable, and each ideal $(\pi,x_i)$ becomes principal on this
model; the stated ambient blowups restrict to isomorphisms there.

At a point where $\bar Q=0$, transversality supplies one additional
relative parameter $y$ and, after absorbing a unit, the actual etale
equation is
\[
                         x_1\cdots x_s=\pi y.
 \tag{MX21.5}\label{r21mx:chart0}
\]
The relative dimension is four, so at most five displayed parameters
are needed.  At the first blowup in $(\pi,x_1)$, the saturated
$x_1$-chart has $\pi=x_1t_1$ and
$x_2\cdots x_s=t_1y$; the other chart has $x_1=\pi u_1$ and
$y=u_1x_2\cdots x_s$.  Repeating the first chart with the next
ideal $(t_1,x_2)$, and continuing, gives the saturated charts
\[
 \pi=x_1\cdots x_{a-1}t_{a-1},\qquad
 x_a=t_{a-1}u_a,\qquad
 y=u_ax_{a+1}\cdots x_s.
 \tag{MX21.6}\label{r21mx:charts}
\]
The final continuing chart is strict semistable as well.  Eliminating
$x_a,y$ leaves regular parameters with $\pi$ a square-free product.
These charts prove regularity and reduced strict semistability and
show that the vertical components are precisely the old factor
components.  Projectivity follows from the actual ambient blowups.
The closure and its strict transform are integral, since the smooth
geometrically integral generic fiber is dense, and they are flat over
the trait.  The local charts also verify flatness directly.

For completeness the stratum calculation in these same charts has
the following global description.  The component belonging to an
earlier factor is modified by a later factor at the ideal
$(\bar f_j,\bar Q)$ on that component; the later component is
unchanged by this step.  On an intersection of factors indexed by
$I$, a factor already in $I$ gives a Cartier ideal and has no effect.
The remaining nontrivial ideals are exactly those with $j>\max(I)$,
and the induced order on this intersection is the reverse order in
\eqref{r21mx:centers}.  The two affine Rees charts of each such ideal
are the corresponding restrictions of \eqref{r21mx:charts}.

These local ideals glue: replacing any defining section by a local
unit multiple does not change the ideal of its zero divisor together
with the zero divisor of $\bar Q$.  Equality of the line bundles of
different factors is not used.  Every $Z_{I,j}$ is smooth of the
asserted codimension by the given transverse-cut hypotheses.  A
preceding center induces on a subsequent $Z_{I,j}$ the ideal of the
remaining factor, because $\bar Q$ already vanishes there.  This is
Cartier, so its blowup is the identity.  On zero-dimensional centers
the additional intersections are empty: they would require four
distinct outer hyperplanes.  This proves the assertion about strict
transforms and completes the actual stratum calculation.
\end{proof}

\subsection{Full component positivity and the native double gate}

\begin{lemma}[The full middle form on the resulting components]
\label{r21mx:positive}
For every component $D_i$, its native numerical middle intersection
form on $N^2(D_i)_{\mathbf R}$ is nondegenerate with negative index
$\rho(D_i)-1$.  For every ample class $H_i$ its primitive part
$\ker(H_i:N^2(D_i)\to N^3(D_i))$ is positive definite.
\end{lemma}

\begin{proof}
On $S\times\mathbf P^2$, write $h$ for the projective-factor class.
The actual projective-bundle formula, followed by numerical duality,
gives the orthogonal decomposition
\[
 N^2(S\times\mathbf P^2)
    =\bigl(N^2(S)\oplus\mathbf Qh^2\bigr)\ \perp\ N^1(S)h.
 \tag{MX21.7}\label{r21mx:product-form}
\]
The first summand is a hyperbolic plane, and the form on the second
is the surface divisor intersection form, of signature
$(1,\rho(S)-1)$.  The total negative index is $\rho(S)$, equal to
$\rho(S\times\mathbf P^2)-1$.  On $C\times\mathbf P^3$ the middle
group has basis $[\mathrm{point}]h,h^2$ after choosing a rational
degree-one zero-cycle class; it is one hyperbolic plane and its
negative index is $1=\rho(C\times\mathbf P^3)-1$.
Native rational zero-cycle degrees justify these statements even
without a native rational point.

If $f:\widetilde D=\operatorname{Bl}_ZD\to D$ blows up a smooth
surface in a smooth fourfold, with exceptional diagram
$j:E\hookrightarrow\widetilde D$, $p:E\to Z$, the actual Chow
decomposition is
\[
 \operatorname{CH}^2(\widetilde D)_{\mathbf Q}
    =f^*\operatorname{CH}^2(D)_{\mathbf Q}
       \oplus j_*p^*\operatorname{CH}^1(Z)_{\mathbf Q}.
 \tag{MX21.8}\label{r21mx:blowup}
\]
The two summands are orthogonal: the putative mixed pairing is a
codimension-three intersection on the surface $Z$.  The exceptional
pairing is
\[
     \deg(j_*p^*a\cdot j_*p^*b)=-\deg_Z(ab),
 \tag{MX21.9}\label{r21mx:exceptional}
\]
by $j^*j_*=-\xi$ and $p_*\xi=1$.
The numerical decomposition is therefore the same direct sum, since
both summand pairings are nondegenerate.  Surface divisor Hodge index
shows that a connected center adds exactly one negative direction,
and the Picard rank increases by one.  Disjoint connected centers
are treated componentwise, retaining their native degree pairings.
Induction through \eqref{r21mx:centers} proves the stated index.

For an arbitrary ample $H_i$, divisor Hodge index makes the pairing
on $H_iN^1(D_i)$ have signature $(1,\rho(D_i)-1)$, and makes
multiplication by $H_i$ on $N^1$ injective.  Its orthogonal complement
is precisely the displayed primitive middle group, by numerical
duality.  Since all negative directions have already occurred in
$H_iN^1$, that complement is positive definite.  This last argument
permits the restrictions of any one genuinely relatively ample
line bundle on the projective model; a pulled-back nef class is not
silently used as ample.
\end{proof}

\begin{lemma}[Curve comparison on every native double]
\label{r21mx:double}
Every double stratum of $\mathcal X$ satisfies the frozen CD17
assertion that an actual native rational curve numerically equivalent
to zero has zero geometric $\ell$-adic class.  This assertion holds
anew after every finite residue-field extension.
\end{lemma}

\begin{proof}
The original doubles are $S\times\mathbf P^1$ and
$C\times\mathbf P^2$.  CD17's projective-bundle comparison applies
to both: the only base numerical comparisons needed are surface
divisors and zero-cycle degree, or curve zero-cycle degree.  It does
not require divisor Tate for $S$.  By
Proposition~\ref{r21mx:resolution} every transformed double is an
iterated blowup in the actual smooth curves $Z_{I,j}$.  The frozen
CD17 smooth-blowup comparison applies to each such curve center.
The same argument includes all new cycles after a finite residue
extension; no descent of those cycles is assumed.
\end{proof}

\begin{theorem}[Complete local Green domain for the mixed smoothing]
\label{r21mx:green}
For every datum \eqref{r21mx:central}--\eqref{r21mx:generic}, the
actual model above has $U_2=0$.  Its generic fourfold has full strong G
in every codimension.  After every finite valued extension $L/K$ the
same complete assertion holds for every new specified geometrically
homologically trivial rational cycle on $X_L$.
\end{theorem}

\begin{proof}
Choose one relatively ample line bundle on $\mathcal X$.  PU18
identifies the full $U_2$ with its unique primitive representative
subspace in $\ker\rho_2\cap\operatorname{im}\sigma_2$ and proves that
subspace totally isotropic.  By Lemma~\ref{r21mx:positive} the ambient
primitive middle form is positive definite.  Hence the subspace is
zero.  Equivalently apply FP21 with the identity projector on every
component: its stronger full-$L_2$ confinement is then automatic.
This retains the genuine triple and quadruple strata; it does not
assert $\ker A=\ker\rho$ or $K_{\rm comp}=\ker A$.

Lemma~\ref{r21mx:double} closes the CD17/CL17 representative gate.
The frozen exact detector, or the MZ20 consequence of $U_2=0$,
then gives codimension-two strong G for every actual original
homological Chow class.  The frozen divisor, one-cycle and zero-cycle
theorems give all other fourfold codimensions.

Every finite unramified extension preserves the smooth strata and
the ordered charts.  Reapply the complete component positivity and
double-comparison arguments there to its entire new Chow groups.
The model is reduced strict semistable, so the frozen RI18
Corollary \texttt{r18ri:all-finite}, from
\texttt{inertial\_triviality\_and\_ramified\_norm\_detection.tex},
propagates the full theorem from all unramified fields to every finite
valued extension.  Its ramified norm argument and final finite trace
retain the complete field degrees, including wild and residue factors.
No new cycle over $L$ is assumed to be pulled back from $K$.

Finally VT17 returns, for each specified final cycle $Z$, an ordinary
actual Green form on exactly $X_L^{\rm an}\setminus|Z|^{\rm an}$,
with its identity against all actual compact complementary tests and
the extending original absolute value.  This is a class-kernel
argument followed by exact-support reconstruction, not an arbitrary
proper actual pushforward.
\end{proof}

\subsection{A full weight-row calculation excluding potential good reduction}

Let
\[
 A_\ell^1(S)=\operatorname{im}\left(
 \operatorname{NS}(S_{\bar k})\otimes\mathbf Q_\ell(-1)
          \xrightarrow{c_1}H^2(S_{\bar k},\mathbf Q_\ell)\right),
 \qquad T=A_\ell^1(S)^\perp,
 \tag{MX21.10}\label{r21mx:T}
\]
where the orthogonal complement is taken in $H^2(S_{\bar k},\mathbf Q_\ell)$
for its degree pairing.  The Tate twist in the divisor subspace is
displayed to retain its Frobenius weight two.  Its intersection form
is nondegenerate by divisor Hodge
index and the divisor numerical/homological comparison, so this is a
direct complement.  Both summands are Frobenius stable.  The statement
does not assume that every Frobenius-invariant class is algebraic.

\begin{theorem}[A nonzero weight-two subquotient]
\label{r21mx:bad}
If $T\ne0$, the weight spectral sequence for the actual model
$\mathcal X$ has a Frobenius-stable direct summand
\[
      T\otimes_{\mathbf Q_\ell}\mathbf Q_\ell^{\binom{d-1}{3}}
                             \subset E_2^{2,2}.
 \tag{MX21.11}\label{r21mx:weight2}
\]
It survives in the weight-two graded piece of
$H^4(X_{\bar K},\mathbf Q_\ell)$.  Consequently $X$ has no potential
good reduction.  No general weight--monodromy conjecture, and no
claim that the entire $E_2^{2,2}$ equals this summand, is used.
\end{theorem}

\begin{proof}
Use the Rapoport--Zink convention already frozen in HM18:
\[
 E_1^{a,b}=\bigoplus_{i\ge\max(0,-a)}
 H^{b-2i}(D^{[a+2i+1]}_{\bar k},\mathbf Q_\ell)(-i)
       \Longrightarrow H^{a+b}(X_{\bar K},\mathbf Q_\ell).
 \tag{MX21.12}\label{r21mx:rz}
\]
Here $D^{[s]}$ is the disjoint union of the actual $s$-fold strata.
The differential is the sum of alternating restrictions and Gysins.
The entire row $b=2$ is
\[
 E_1^{-1,2}=H^0(D^{[2]})(-1),\qquad
 E_1^{a,2}=H^2(D^{[a+1]})\oplus H^0(D^{[a+3]})(-1)
                                      \quad(a\ge0),
 \tag{MX21.13}\label{r21mx:row}
\]
with empty-stratum and impossible-degree terms zero.  There are no
terms for $a<-1$.  In this model $D^{[5]}=\varnothing$, but the
formula retains the Gysin summands in the preceding terms.

Every smooth blowup in \eqref{r21mx:centers} adds to $H^2$ only
$H^0$ of its center with Tate twist $(-1)$, represented by exceptional
divisors.  Thus it adds only algebraic divisor classes.  For a stratum
not containing index $0$, its old $H^2$ is
$H^2(S)\oplus H^2(\mathbf P^{3-|I|})$, and its unique nondivisor
summand is a copy of $T$, pulled back from $S$.  Such strata exist
exactly for $1\le |I|\le3$.  Every stratum containing $0$ starts from
$C\times\mathbf P^{4-|I|}$, whose $H^2$ is generated by divisor
classes, and remains so after the indicated blowups.

The pullback of $T$ to $C$ is zero: its trace on the curve is its
intersection with $[C]$, which is zero by definition of $T$, and
$H^2(C)$ is the degree line.  Therefore a face restriction between
two strata with indices outside $0$ is the identity on these pulled
back copies of $T$, and a face restriction into a stratum containing
$0$ is zero.  Algebraic divisor classes restrict to algebraic
divisor classes.  Every Gysin from the $H^0(-1)$ summands of
\eqref{r21mx:row} also has algebraic divisor image.  Hence projection
onto the copies of $T$ commutes with the entire row differential;
the pullbacks from $S$ give a splitting as a complex.  This checks
both the restriction terms and all extra Gysin terms, rather than
discarding the latter on a dimension guess.

The resulting direct subcomplex is
\[
 T^{\oplus d}\xrightarrow{\delta}
 T^{\oplus\binom d2}\xrightarrow{\delta}
 T^{\oplus\binom d3}\longrightarrow0,
 \tag{MX21.14}\label{r21mx:Tcomplex}
\]
in columns $a=0,1,2$.  It is $T$ tensored with the cochain complex
of the two-skeleton of the $(d-1)$-simplex.  The augmented full-simplex
complex is exact.  Truncation at degree two leaves
\[
 \dim H^2\bigl(\operatorname{skel}_2\Delta^{d-1},\mathbf Q_\ell\bigr)
   =\binom d3-\binom d2+d-1=\binom{d-1}{3}.
 \tag{MX21.15}\label{r21mx:binomial}
\]
This proves the claimed direct summand on $E_2$.

Every term $E_1^{a,b}$ is pure of weight $b$, because the strata are
smooth projective over the finite field.  The higher differentials
$d_r$, $r\ge2$, connect different weights and are zero.  Thus the
summand survives in the abutment filtration of $H^4$ and its
Frobenius eigenvalues have weight two.  This uses finite-field Weil
purity and the weight spectral sequence, not weight--monodromy.

If $X$ acquired good reduction over a finite extension of residue
degree $f$, proper smooth base change and Weil purity would make
every eigenvalue on $H^4$ have complex absolute value $q^{2f}$.
On the above nonzero graded piece the eigenvalues instead have
absolute value $q^f$.  The discrepancy survives every finite
extension: inertia acts trivially on the graded pieces of
\eqref{r21mx:rz}, and a new Frobenius acts there by the $f$th power
of the old Frobenius.  This contradiction excludes potential good
reduction, including after a ramified extension.
\end{proof}

\paragraph{Primary source for the spectral sequence.}
The conventions and differential description are in T.~Saito,
\emph{Weight spectral sequences and independence of $\ell$},
Proposition 2.2.3 and Corollary 2.2.4, author manuscript
\url{https://www.ms.u-tokyo.ac.jp/~t-saito/pp/wmr2r.pdf}.
The finite-field purity input is Deligne's Weil theorem.  This is
the same foundational input as frozen HM18; the coefficient-$T$
row calculation above is supplied explicitly.

\subsection{Nonempty local families and exact limits of the extension}

\begin{proposition}[Existence after a finite unramified preparation]
\label{r21mx:existence}
For every supplied good surface $\mathcal S/R$ and every $d\ge4$,
one may choose $\mathcal A$ sufficiently relatively ample and make a
finite unramified extension so that all the above hypotheses hold.
For fixed prepared central factors, the smooth admissible $Q$ lifts
form a Zariski-dense set in the corresponding affine coefficient
space over the prepared fraction field.  If $T\ne0$, every one of
these admissible smooth fibers has the obstruction-free complete
local Green domain and has no potential good reduction.
\end{proposition}

\begin{proof}
Take a sufficiently large tensor power of a relatively ample line
bundle so that it is relatively very ample and its sections commute
with reduction.  Serre vanishing supplies this last condition.
Over $\bar k$, classical Bertini supplies a smooth geometrically
integral $C$ and the finite hyperplane arrangement in general
position.  The line bundle $\mathcal A\boxtimes\mathcal O(d)$ is
very ample, including on each smooth central stratum.  Its ambient
linear system restricts to a very ample system there.  Simultaneous
Bertini and transversality for the finitely many strata give a
nonempty open set of $\bar Q$.  A general section cuts a curve
stratum in a reduced finite set and a stratum of dimension at least
two in a smooth geometrically integral divisor.

All these finitely many choices are defined over a finite residue
extension, hence over a finite unramified extension of $R$.
The selected $f_0,Q$ lift because of the chosen cohomology-and-base-
change surjectivity; for $Q$ the projective-factor cohomology adds
no obstruction.  Hyperplane forms lift directly.  Fix those residue
choices.  The lifts of $Q$ form an open residue coset in a full
integral coefficient lattice.  Such a coset is Zariski dense over
the infinite fraction field: a polynomial vanishing on it is zero,
by induction in the number of variables.

Over that fraction field the map $Q\mapsto f_0\prod l_j+\pi Q$
is an affine isomorphism of the full section space.  The smooth
geometrically integral divisors of the very ample linear system
form a nonempty Zariski open set.  It therefore meets the residue
coset densely, giving exactly the required generic smooth fibers.
Apply Theorems~\ref{r21mx:green} and \ref{r21mx:bad}.
\end{proof}

For a concrete nonzero-$T$ good surface, take
$\mathcal S=\mathcal E\times\mathcal E$ over $\mathbf Z_5$ with
$\mathcal E:y^2=x^3+x$.  Its special elliptic curve has four points,
Frobenius polynomial $t^2-2t+5$, and the native order-four
automorphism $(x,y)\mapsto(-x,2y)$.  The same elementary CM calculation
used in FP21 gives $\rho(S_{\bar k})=4$, $\dim H^2(S)=6$, so
$\dim T=2$.  One can also see the upper bound directly: the two
remaining eigenvalues are $(1\pm2i)^2$, whose ratio to $5$ is not
a root of unity, while every divisor class has a Tate eigenvalue
after some finite extension.  The factor, diagonal and automorphism
graph divisors supply four independent classes.  With $d=4$,
\eqref{r21mx:weight2} already has dimension two.  The subsequent
Bertini preparation produces actual, potentially-bad-reduction
fourfold fibers; it is not a vertical modification of this good
surface's square.

\paragraph{What is new here, and what remains unproved.}
PG18 and SC18 assume that all product factors are positive powers of
one relatively very ample line bundle.  Here $H_0=C\times\mathbf P^3$
and $H_j=S\times\mathbf P^2$ are nonample in $S\times\mathbf P^3$;
their factor line bundles are $\mathcal A\boxtimes\mathcal O$ and
$\mathcal O\boxtimes\mathcal O(1)$.  Thus the frozen common-ample-factor
theorem does not apply to this presentation.  The proof above supplies
the required model, its full component positivity, and its native
double comparison directly.  The weight-two calculation separately
proves that these fibers need not have potential good reduction.

We do not assert that no member could have another earlier qualifying
presentation or satisfy a more general earlier motivic criterion.
The newly proved class is this explicit local smoothing presentation,
stable for all new cycles after finite extension; a claim of strict
set-theoretic non-overlap with every earlier G class is not made.
The argument does not prove every smooth divisor in
$|\mathcal A_K\boxtimes\mathcal O(d)|$ has G.

For a variety over a number field one may apply the local theorem
at a completion with such supplied data.  A complete BB height on
the full global homological Chow groups still requires the frozen
local hypotheses at every other finite place.  No such global
coefficient or spreading construction is supplied here.  General G
and arbitrary-variety complete unconditional heights remain open.
